% Ad Familiares 4.12 (ad-familiares-04-12) --- English translation
% Translated by Claude (Anthropic) from the Latin (Perseus).
% Style guide: STYLE.md.

\ciceroLetterOpener{SERVIVS CICERONI sal. PLVR.}

\ciceroSection{1} Although I know that the news I shall bring you is not the most welcome, still --- since chance and nature have the mastery over us --- it seemed right that, whatever the case might be, I should inform you. On the tenth day before the Kalends of June, when I had crossed by ship from Epidaurus to the Piraeus, I met there our colleague M. Marcellus, and I spent that day there in order to be with him. On the next day, when I had taken my leave of him, with the intention of going from Athens into Boeotia and finishing out the remainder of my judicial business, he, by his own account, was on the point of sailing past Cape Malea toward Italy.

\ciceroSection{2} On the third day after that, when I had it in mind to set out from Athens, at about the tenth hour of the night P. Postumius, his intimate, came to me and reported that our colleague M. Marcellus had been struck after dinner by P. Magius Cilo, his own intimate, with a dagger, and had received two wounds, one in the belly, the other in the head beside the ear; that he hoped, nonetheless, that he might live; that Magius had killed himself afterward; that he had been sent to me by Marcellus to bring me this news and to ask that I send him doctors. So I assembled doctors and set out at once for the spot, at first light. When I was not far from the Piraeus, the boy Acidim came to meet me with a wax tablet on which was written that a little before dawn Marcellus had met his end.

\ciceroSection{3} So a most illustrious man has been carried off by a most worthless one by the harshest of deaths; and the man whom his enemies had spared for the sake of his standing has found a friend to deal him death. I went on, even so, to his tent. I found two freedmen and a very few slaves; the rest, they said, had fled in terror, because their master had been killed in front of his tent. I was compelled, in the very litter in which I had myself been carried, and with my own bearers, to bring him back into the city; and there, with such resources as Athens afforded, I saw to it that the funeral made for him should be sufficiently splendid. From the Athenians I could not win the granting of a place of burial within the city, since they said they were prevented by religious scruple --- and indeed they had never before granted it to anyone. The next best thing they did permit us: that we should bury him in whichever gymnasium we wished. We chose, in the most famous gymnasium in the world, the place of the Academy; and there we cremated him, and afterwards saw to it that those same Athenians should commission, on that same spot, a marble monument for him. Thus, in what duties were ours to discharge --- in virtue of our shared office and our personal connection --- we have done him every service, in life and in death. Farewell. Dispatched the day before the Kalends of June, from Athens.
